\subsection{Temporal Stability of Marginal Predictions}
\label{subsec:blocked_marginal_stability}

The marginal contribution defined in the preceding section is estimated over the
complete validation period. A positive value may nevertheless be concentrated in
a short historical interval. The robustness objective therefore examines the
distribution of the candidate's marginal contribution across contiguous validation
subperiods.

This calculation reuses the two combined-model predictions already fitted for
the marginal-value objective. In particular, the robustness calculation does not
refit either model. This distinction is deliberate. The combined model is used as
a flexible diagnostic of the conditional information contained in the candidate,
rather than as the final forecasting model of the investment strategy. Repeatedly
refitting it would combine variation in the factor's contribution with variation
in the model-estimation procedure. The present objective instead asks whether the
contribution identified by one fixed in-sample fit is distributed consistently
through the validation period.

Let
\[
    \widehat R^{(h)}_{\mathcal B\cup\{X\},t,a}
    =
    \widehat f_{\mathcal B\cup\{X\}}
    \left(
        \mathbf z^{\mathcal B}_{t,a},z^X_{t,a}
    \right)
\]
and
\[
    \widehat R^{(h)}_{\mathcal B,t,a}
    =
    \widehat f_{\mathcal B}
    \left(
        \mathbf z^{\mathcal B}_{t,a}
    \right)
\]
denote the two prediction panels fitted on the in-sample period \(\mathcal I\)
for the marginal-value calculation. The standardisation parameters, model
coefficients and tree structures are therefore identical in the marginal-value
and robustness objectives.

Validation observations whose forward-return label would require a price from the
final test period are removed. Define the remaining validation index by
\[
    \mathcal V_h
    =
    \left\{
        t\in\mathcal V:t+h\in\mathcal I\cup\mathcal V
    \right\}.
\]
The ordered index \(\mathcal V_h\) is partitioned into \(K\) contiguous,
non-overlapping blocks of approximately equal length,
\[
    \mathcal V_h
    =
    \mathcal V_1\mathbin{\dot\cup}\cdots
    \mathbin{\dot\cup}\mathcal V_K.
\]
The baseline implementation uses \(K=4\). Contiguous blocks preserve the
temporal structure of the data and allow the objective to detect contributions
that are confined to one part of the validation period.

For block \(k\), define the candidate's fixed-prediction marginal contribution by
\begin{equation}
    \Delta IC_k
    =
    \operatorname{IC}_{\mathcal V_k}
    \left(
        \widehat R^{(h)}_{\mathcal B\cup\{X\}}
    \right)
    -
    \operatorname{IC}_{\mathcal V_k}
    \left(
        \widehat R^{(h)}_{\mathcal B}
    \right).
    \label{eq:block_marginal_ic}
\end{equation}
If the current factor book is empty, the second term in
Equation~\eqref{eq:block_marginal_ic} is defined as zero. Only blocks for which
the relevant information coefficients are defined enter the aggregation. For
the resulting set of valid blocks \(\mathcal K_X\), let
\[
    \overline{\Delta IC}_X
    =
    \frac{1}{|\mathcal K_X|}
    \sum_{k\in\mathcal K_X}\Delta IC_k
\]
and
\[
    \sigma_{\Delta IC,X}
    =
    \left[
        \frac{1}{|\mathcal K_X|}
        \sum_{k\in\mathcal K_X}
        \left(
            \Delta IC_k-\overline{\Delta IC}_X
        \right)^2
    \right]^{1/2}.
\]
The basic temporal-stability score is
\begin{equation}
    F_{\mathrm{block}}
    \left(
        X\mid\mathcal B
    \right)
    =
    \overline{\Delta IC}_X
    -
    \lambda_{\mathrm{std}}\sigma_{\Delta IC,X},
    \label{eq:block_stability_score}
\end{equation}
where \(\lambda_{\mathrm{std}}\geq0\) determines the penalty assigned to temporal
dispersion. The implementation uses \(\lambda_{\mathrm{std}}=1\). A candidate
therefore receives a high value only when its average marginal contribution is
positive and does not depend strongly on the choice of validation subperiod.
Unlike a sign-aligned stability statistic, Equation~\eqref{eq:block_stability_score}
does not reward a candidate for contributing negatively in a consistent manner.

For diagnostic purposes, the evaluator also computes the raw factor IC separately
on every block,
\[
    IC_k(X)
    =
    \operatorname{IC}_{\mathcal V_k}(X).
\]
The mean and standard deviation of these values describe the temporal stability
of the candidate's unconditional association with returns. They do not enter the
selection objective. This separation is necessary for conditioning variables,
such as volatility states, whose standalone IC may remain close to zero even when
their nonlinear marginal contribution to the factor book is positive.

The parameter-plateau penalty from the preceding construction is retained. Let
\(X^{(1)},\ldots,X^{(J)}\) denote deterministic variants obtained by perturbing
recognised rolling-window parameters, and define
\[
    s_X
    =
    \begin{cases}
        1, & IC_{\mathcal V}(X)\geq0,\\
        -1, & IC_{\mathcal V}(X)<0.
    \end{cases}
\]
The penalty is
\begin{equation}
    P_{\mathrm{plat}}(X)
    =
    \max\left\{
        0,
        s_XIC_{\mathcal V}(X)
        -
        \frac{1}{J}
        \sum_{j=1}^{J}
        s_XIC_{\mathcal V}\left(X^{(j)}\right)
    \right\}.
    \label{eq:block_plateau_penalty}
\end{equation}
The baseline implementation perturbs eligible windows by approximately
\(\pm10\%\). Candidates without recognised window parameters receive no plateau
penalty.

If the researcher declares an expected sign \(e_X\in\{-1,1\}\) before
evaluation, the hypothesis-sign term is
\[
    B_{\mathrm{sign}}(X)
    =
    \begin{cases}
        b,
        & \operatorname{sign}\left(IC_{\mathcal V}(X)\right)=e_X,\\
        -b,
        & \operatorname{sign}\left(IC_{\mathcal V}(X)\right)\neq e_X,\\
        0,
        & e_X\text{ or }IC_{\mathcal V}(X)\text{ is undefined},
    \end{cases}
\]
with \(b=0.02\). An optional signal-perturbation penalty
\(P_{\xi}(X)\) can be included as an ablation. Its baseline weight is zero.

The complete robustness objective is
\begin{equation}
\begin{split}
    F_{\mathrm{rob}}
    \left(
        X\mid\mathcal B
    \right)
    ={}&
    \overline{\Delta IC}_X
    -
    \lambda_{\mathrm{std}}\sigma_{\Delta IC,X}
    -
    w_{\mathrm{plat}}P_{\mathrm{plat}}(X)\\
    &-
    w_{\xi}P_{\xi}(X)
    +
    B_{\mathrm{sign}}(X).
\end{split}
\label{eq:complete_block_robustness}
\end{equation}
The baseline values are
\[
    \lambda_{\mathrm{std}}=1,
    \qquad
    w_{\mathrm{plat}}=1,
    \qquad
    w_{\xi}=0.
\]

This objective is an internal development-time measure. Since the evolutionary
researcher repeatedly observes validation-based fitness, the validation period is
eventually adapted to by the search. Equation~\eqref{eq:complete_block_robustness}
therefore does not replace the final untouched test and walk-forward evaluations.
Those procedures assess the generalisation of the complete research process,
whereas the present objective measures whether one IS-fitted marginal contribution
is temporally dispersed within the validation sample.
